% Page 42

\lettrine{U}{n} autre membre de cette famille, Jean de Moret, qui est bénéficier de Ste Cécile d'Albi (quelque chose comme chanoine titulaire), est bénéficiaire de la « Rectorie du Puy Saint Georges, près de Valdériès, et il l'arrente le 16 juin 1529 pour le compte de François de Lautrec pour la somme de 400 livres. On le trouve encore à Sainte Cécile en 1540. Un Jacques de Moret nous est signalé le 11 janvier 1567 dans un accord passé avec Antoine Laveygne. Pierre de Mammesse (orthographe incertaine) se reconnaît débiteur le 9 mai 1580 à l'égard de Pierre de Moret. Le débiteur est un notable puisque lieutenant principal dans le Castrais. La même année, un Antoine Caussé achète un cheval pour 55 écus et Louis de Moret, le soi-disant fondateur de la Chapellenie. Un autre de Moret, un peu plus tard en 1604 et le 3 juillet, se trouve être le prévôt de St Salvy d'Albi et il prétend devoir jouir d'une rente de 50 écus d'or à prendre sur les revenus du prieuré de Cambon et de Montels (son annexe) et il exhibe par devant notaire une lettre du Pape Clément VIII qui le fait bénéficier de cette rente....

\medskip

Quand, en 1945, fut révisé notre cadastre, on eut pour ce faire, des administrateurs ignorant notre patois languedocien. C'est un peu regrettable, car ne connaissant pas l'occitan, ils massacrèrent quelque peu les noms, les vocables et ils appelèrent le « rec de la Cappélanyé » « ruisseau de la Chapelaine ». Ceci sent le roman -- ce n'est pas très méchant et c'est bien savoureux, mais c'est : « Ruisseau de la CHAPELENIE ».
